% ============================================================
% Section 3.2 — P(F) and Continuous Topology Change via Soft Cluster Weights
% Target: ~1.0–1.2 page of 2-column IEEE conference format
% Status: 3rd DRAFT 2026-05-15 (REORDERED — was the Jacobi PBD section.
%         Now P(F) lives here AND is followed by the explicit
%         w_i := α bridge (Eq.~\ref{eq:bridge}) + Algorithm 1 + the
%         cycle-scoped buffer pattern. Bridge is the climax, not a
%         forward reference. Filename kept; content rebuilt.)
% ============================================================

\subsection{$P(F)$ and Continuous Topology Change via Soft Cluster Weights}
\label{sec:pf-cutting}

% --- Block 1: high-level definition + edge-level instantiation justifications ---
We define $P(F)$ as a differentiable score for the survival of a
topological feature $F$ under intruder geometry. Although the symbol
$F$ may suggest a $2$-simplex (face) reading, we adopt the edge level
($1$-simplex) as the primitive representation because edge-level
scores factorize naturally into face- and tet-level scores when
needed, without separately tracking those higher-order variables. Let
$F$ denote the event ``the feature is intact'', let $e$, $f$, and $T$
denote a single edge, face, and tet respectively, and let $e \in f$
(resp. $e \in T$) range over the bounding edges of $f$ (resp. $T$).
Treating the survival of a face or tet as the joint event ``all of
its constituent edges are intact'', and assuming the conditional
independence of those edge events given the intruder $I$, we obtain
the modeling approximation
\begin{equation}
  P(F \mid f, I) \!=\! \prod_{e \in f}\!P(F \mid e, I),
  \quad
  P(F \mid T, I) \!=\! \prod_{e \in T}\!P(F \mid e, I).
  \label{eq:simplicial-product}
\end{equation}
The edge representation is therefore sufficient to induce face- and
tet-level survival scores via this boundary-edge factorization,
without explicitly maintaining separate higher-simplex cut variables.

% --- Block 2: P(F) precise definition + properties ---
We now make $P(F)$ precise. Consider a tetrahedral mesh
$\mathcal{M} = (\mathcal{V}, \mathcal{E}, \mathcal{F}, \mathcal{T})$
with current vertex positions $\mathbf{p}: \mathcal{V} \to \mathbb{R}^3$
and an intruder $I = (\mathbf{c}_I, \rho_I)$ defined by its center
$\mathbf{c}_I$ and radius $\rho_I \!\geq\! 0$. The intruder is
instantiated either as the cutting tool (a sphere of physical radius
$\rho_I \!>\! 0$) or, at mesh initialization, as the nearest
non-adjacent vertex of each edge ($\rho_I \!=\! 0$, a point proxy
that yields a mesh-density-aware baseline $P(F)_e^{\text{rest}}$).
For an edge $e = (v_a, v_b)$ with midpoint
$\mathbf{m}_e = \tfrac{1}{2}(\mathbf{p}_a + \mathbf{p}_b)$ and
half-length $r_e = \tfrac{1}{2}\|\mathbf{p}_a - \mathbf{p}_b\|$, the
signed distance $d$ between the edge's circumsphere
$B(\mathbf{m}_e, r_e)$ and the intruder $B(\mathbf{c}_I, \rho_I)$ is
\begin{equation}
  d(e, I) \;=\; \|\mathbf{m}_e - \mathbf{c}_I\| - \rho_I - r_e ,
  \label{eq:sdist}
\end{equation}
positive when the spheres are disjoint, zero at tangency, and negative
under penetration. The Euclidean norm $\|\mathbf{m}_e - \mathbf{c}_I\|$
is regularized as $\sqrt{\|\mathbf{m}_e - \mathbf{c}_I\|^2 + \varepsilon^2}$
to remove its singularity at $\mathbf{m}_e \!=\! \mathbf{c}_I$. The differentiable existence probability is then
\begin{equation}
  P(F \mid e, I)
  \;=\; \sigma\!\left(\kappa\,\frac{d(e, I)}{r_e}\right) ,
  \label{eq:pf}
\end{equation}
where $\sigma(x) = \frac{1}{1+e^{-x}}$ and $\kappa > 0$ is a sharpness
parameter (distinct from the bulk stiffness $k$ introduced in
\S\ref{sec:sysid-framework}). The dimensionless ratio $|d|/r_e$ is
scale-invariant in the mesh, so $\kappa$ carries a unit-free meaning.
By sigmoid saturation, $P(F) \to 1$ when the intruder is far
($d \gg 0$), $P(F) \to 0$ under deep penetration ($d \ll 0$), and
$P(F) = \tfrac{1}{2}$ at tangency. The construction is $C^\infty$ in
vertex and intruder positions wherever $r_e > 0$; the only singularity
at $r_e = 0$ is regularized in practice as
$r_e \!\leftarrow\! \sqrt{r_e^2 + \varepsilon^2}$ with
$\varepsilon \!=\! 10^{-4}$, ensuring autograd compatibility throughout.

The baseline $P(F)_e^{\text{rest}}$ used in (\ref{eq:edge-damage})
below is obtained by evaluating (\ref{eq:pf}) once at mesh
initialization with the point-intruder configuration, capturing each
edge's circumsphere occupancy under natural mesh density. This
1-simplex construction parallels the per-face minimum-ball score of
DMesh++~\cite{son2025dmesh2}, restricted to local 1-ring neighbors
for $O(\deg)$ cost per edge.

% --- Block 3: damage chain (P(F) → δ_e → δ_v → α) ---
$P(F)$ propagates to a per-vertex participation mask $\alpha$ through
three smooth maps. First, $P(F)$ drives a per-edge damage state via a
smooth approximation of the positive increment, encoding the physical
irreversibility of cutting:
\begin{equation}
  \delta_e^{\text{next}} \!=\! \delta_e^{\text{prev}} +
    \mathrm{softplus}_{\!+}\!\left(T_e - \delta_e^{\text{prev}}\right),
  \quad
  T_e := \mathrm{softplus}_{\!+}\!\left(1 - \frac{P(F)_e}{P(F)_e^{\text{rest}}}\right),
  \label{eq:edge-damage}
\end{equation}
where $\mathrm{softplus}_{\!+}(x) \!=\! \tfrac{1}{k_s}\log(1 + e^{k_s x})$
smoothly approximates $\max(x, 0)$ (with sharpness $k_s$).
The outer softplus is a differentiable proxy for the positive
increment $[T_e - \delta_e^{\text{prev}}]_+$, so that
$\delta_e^{\text{next}} \!\approx\! \max(\delta_e^{\text{prev}}, T_e)$
while preserving monotonic accumulation. Because the smooth
approximation carries a small positive tail
($\mathrm{softplus}_{\!+}(0) \!=\! \log 2 / k_s$), $\delta_e$ may
exhibit a slow upward drift even when $T_e \!<\! \delta_e^{\text{prev}}$;
the drift rate is bounded by $\log 2 / k_s$ per substep and is
negligible at the cutting horizons used in our experiments
(\S\ref{sec:gradient-verify}).
Second, per-edge damages aggregate to per-vertex damages by averaging
over incident edges, under the same smooth clamp:
\begin{equation}
  \delta_v^{\text{next}} \!=\! \delta_v^{\text{prev}} +
    \mathrm{softplus}_{\!+}\!\left(\bar{\delta}_e^{(v)} -
       \delta_v^{\text{prev}}\right),
  \quad
  \bar{\delta}_e^{(v)} := \frac{1}{|E(v)|}
      \sum_{e \in E(v)} \delta_e^{\text{next}},
  \label{eq:vertex-damage}
\end{equation}
where $E(v) = \{e \in \mathcal{E} : v \in e\}$ is the set of edges
incident to $v$. The mean (rather than sum) aggregator ensures the
threshold $\tau$ below has a structure-invariant meaning: interior
vertices ($|E(v)| \!\sim\! 6$) and boundary vertices
($|E(v)| \!\sim\! 2$) reach the cut regime under comparable per-edge
damage. Third, the vertex damage drives the sigmoid participation
mask
\begin{equation}
  \alpha_v \!=\! \sigma\!\left(-s\,(\delta_v - \tau)\right) \in [0, 1],
  \label{eq:soft-removal}
\end{equation}
with sharpness $s$ and threshold $\tau$. $\alpha_v \!=\! 1$ denotes
full participation in the dynamics; $\alpha_v \!=\! 0$ denotes
dynamic removal; intermediate values interpolate smoothly between the
two.

% --- Block 4: BRIDGE — w_i := α (the climax of the section) ---
The two layers of the simulator---continuous topology
((\ref{eq:pf})--(\ref{eq:soft-removal})) and shape-matching dynamics
((\ref{eq:centroid})--(\ref{eq:cluster-correction}))---are coupled by
the single assignment
\begin{equation}
  w_i \;:=\; \alpha_{v_i}.
  \label{eq:bridge}
\end{equation}
The topology and dynamics layers are coupled only through the
participation weight (\ref{eq:bridge}). The same $\alpha_{v_i}$
enters the centroid (\ref{eq:centroid}), the cross-covariance
(\ref{eq:covariance}), and the per-particle correction
(\ref{eq:cluster-correction}) with no auxiliary thresholds or discrete
switches: $\alpha_{v_i} \!=\! 1$ recovers the standard shape-matching
solve; $\alpha_{v_i} \!=\! 0$ removes that vertex's contribution to
every overlapping cluster simultaneously; intermediate values provide
a differentiable relaxation of discrete weight-flipping.

\textbf{Scope of ``continuous topology change.''} Cutting never
modifies $\mathcal{V}, \mathcal{E}, \mathcal{F}, \mathcal{T}$ or the
overlapping-cluster construction, both fixed at initialization;
material loss is encoded entirely in the continuous $\alpha$ field.
This is an \emph{attenuation-based} approximation of topology change
rather than a kinematic separation: when a cut traverses a cluster,
the cut-line vertices have $\alpha \!\to\! 0$ and their reduced
participation softens the cluster's effective coupling, but two
$\alpha \!\approx\! 1$ subregions of the same cluster on opposite
sides of the cut remain bound by a single $R_C, \boldsymbol{\mu}_C$
and cannot detach kinematically. Our experimental cluster scale
(BFS-2 neighborhood, mean $\sim\!30$ vertices, mean edge length
$\sim\!1.5$\,mm) is therefore chosen finer than typical cut depths
($\geq\!5$\,mm) so that most cuts span multiple clusters, with
$\alpha$-driven softening dominating within each cluster and effective
decoupling occurring at cluster boundaries; we report empirical
cavity-formation behaviour in \S\ref{sec:sysid-synth}. Genuine
fragment-detachment with independent rigid-body dynamics for severed
chunks is outside the scope of this formulation, which targets
sys-ID against quasi-static cavity geometry rather than transient
fragment trajectories.

This design preserves the fixed computation graph required by
reverse-mode autodiff while providing a continuously parameterized
participation mask that supports gradient-based identification of the
parameters governing $\alpha$. Algorithm~\ref{alg:ssa-jacobi} summarises
one substep of the integrated solver; Figure~\ref{fig:cluster-jacobi}
illustrates it visually.

% --- Algorithm 1 ---
\begin{algorithm}[!htbp]
\small
\caption{Differentiable PBD substep with $P(F)$ cutting and cycle-scoped buffers.}
\label{alg:ssa-jacobi}
\Input{$\mathbf{p}^{\text{prev}}, \mathbf{v}^{\text{prev}},
       \alpha^{\text{prev}},
       \delta^{\text{e,prev}}, \delta^{\text{v,prev}},
       (\mathbf{c}_I, \rho_I),\, dt,\, K;$\;
       parameters $\kappa, s, \tau, k$ (softplus sharpness $k_s$ fixed)}
\Output{$\mathbf{p}^{\text{next}}, \mathbf{v}^{\text{next}},
        \alpha^{\text{next}},
        \delta^{\text{e,next}}, \delta^{\text{v,next}}$}

$\mathbf{p}^{(0)} \leftarrow \mathbf{p}^{\text{prev}} + dt \cdot \mathbf{v}^{\text{prev}}$
\mycmt{predictor step}\\

\For{$e \in \mathcal{E}$ \KwSty{parallel}}{
  $P(F)_e \leftarrow \sigma\!\left(\kappa\,\frac{d(e, \mathbf{c}_I)}{r_e}\right)$
  \mycmt{Eq.~\ref{eq:pf}}\\

  $T_e \leftarrow
   \mathrm{softplus}_{+}\!\left(1 - \tfrac{P(F)_e}{P(F)_e^{\text{rest}}}\right)$
  \mycmt{smooth target damage}\\

  $\delta^{\text{e,next}}_e \leftarrow \delta^{\text{e,prev}}_e +
   \mathrm{softplus}_{+}\!\left(T_e - \delta^{\text{e,prev}}_e\right)$
  \mycmt{Eq.~\ref{eq:edge-damage}, fresh buffer}\\
}

\For{$v \in \mathcal{V}$ \KwSty{parallel}}{
  $\bar{\delta}^{(v)}_e \leftarrow
   \tfrac{1}{|E(v)|}\sum_{e \in E(v)} \delta^{\text{e,next}}_e$
  \mycmt{mean over incident edges}\\

  $\delta^{\text{v,next}}_v \leftarrow \delta^{\text{v,prev}}_v +
   \mathrm{softplus}_{+}\!\left(\bar{\delta}^{(v)}_e - \delta^{\text{v,prev}}_v\right)$
  \mycmt{Eq.~\ref{eq:vertex-damage}}\\

  $\alpha^{\text{next}}_v \leftarrow
   \sigma\!\left(-s\,(\delta^{\text{v,next}}_v - \tau)\right)$\\
}

\For{$c \leftarrow 0$ \KwTo $K{-}1$}{
  \For{cluster $C$ \KwSty{parallel}}{
    $w_i \leftarrow \alpha^{\text{next}}_i$
    \mycmt{bridge, Eq.~\ref{eq:bridge}}\\

    $\mu \leftarrow \frac{1}{\sum_{i \in C} w_i}\sum_{i \in C} w_i\,\mathbf{p}^{(c)}_i$\\

    $A \leftarrow \sum_{i \in C} w_i\,(\mathbf{p}^{(c)}_i - \mu)
                  (\bar{\mathbf{p}}_i - \bar{\boldsymbol{\mu}})^{\!\top}
                  + \varepsilon I$
    \mycmt{$\varepsilon \!=\! 10^{-6}$ Tikhonov}\\

    $R \leftarrow \mathrm{polar}(A)$\\

    $\Delta_i^{(C)} \leftarrow
     w_i\!\left(R\,(\bar{\mathbf{p}}_i - \bar{\boldsymbol{\mu}})
     + \mu - \mathbf{p}^{(c)}_i\right)$\\
  }
  $\mathbf{p}^{(c+1)}_i \leftarrow \mathbf{p}^{(c)}_i +
   \dfrac{k}{|\mathcal{C}(i)|}\!\sum_{C \in \mathcal{C}(i)} \Delta_i^{(C)}$
  \mycmt{apply with bulk stiffness $k$; mean over clusters containing $i$; fresh buffer}\\
}

$\mathbf{p}^{\text{next}} \leftarrow \mathbf{p}^{(K)},\;\;
 \mathbf{v}^{\text{next}} \leftarrow
 \frac{\mathbf{p}^{\text{next}} - \mathbf{p}^{\text{prev}}}{dt}$\\

\Return $\mathbf{p}^{\text{next}}, \mathbf{v}^{\text{next}},
        \alpha^{\text{next}},
        \delta^{\text{e,next}}, \delta^{\text{v,next}}$
\vspace{6pt}
\end{algorithm}

% --- Block 5: cycle-scoped buffer pattern (engineering discipline) ---
In practice, making Algorithm~\ref{alg:ssa-jacobi} differentiable
requires one implementation discipline: every state buffer is written
to a fresh allocation rather than modified in place. With $K$ Jacobi
iterations, this means $K{+}1$ position buffers; iteration $c$ reads
$\mathbf{p}^{(c)}$ and writes $\mathbf{p}^{(c+1)}$; the same pattern
applies to $\alpha$, $\delta^{\text{e}}$, $\delta^{\text{v}}$, and the
per-iteration cluster rotations. This is a standard prerequisite for
reverse-mode autodiff: in-place writes to tensors whose pre-update
values are needed for backward either raise tape errors or silently
zero out gradients~\cite{difftaichi}. With correct buffers the backward
pass $\partial \mathcal{L}/\partial \theta$ is automatic; we verify in
\S\ref{sec:gradient-verify} that the resulting gradients match finite
differences to machine precision.
